\section{Editor}

\begin{enumerate}
	\item Though this piece came to me as a methods paper, I agree with R. 2 that the best approach is probably to focus foremost on the substantive motivation, lay out your original theoretical perspective, show how that determines the choice of methodology, and explain what you found. That dovetails with R. 3's good suggestions about approaching the presentation, and with R. 1's critiques about the clarity of the central aim of the paper. The key questions are boiled down succinctly by R. 3 who asks that you clarify, ``what are the consequences of these limitations [that you identify in previous approaches]? What don't we understand that, at the end of this article, we will understand better?''
	\begin{itemize}
		\item \textcolor{blue}{ \emph{
		We clarify the purpose of the paper, principally in 
		}}
	\end{itemize}	
	\item In addition, though you certainly sought to avoid wading into the latent space ``vs.'' (T)ERGM ``debate'', I agree with R. 1 that more discussion of your preference for one over the other is called for. To me the most straightforward (if labor intensive), and most transparent way to handle this is to consider TERGM as a robustness check in an appendix.
	\begin{itemize}
		\item \textcolor{blue}{ \emph{
		Insert great response.
		}}
	\end{itemize}		
\end{enumerate}